\section{Justificativa}

Apesar do crescente número de publicações na área de controle de próteses por eletromiografia de superfície (sEMG, do inglês \textit{Surface Electromyography}), ainda existe uma lacuna significativa entre os protótipos de laboratório e dispositivos funcionais, acessíveis e de fácil reprodução. A maioria dos sistemas disponíveis comercialmente apresenta custo elevado, o que restringe seu acesso. Além disso, soluções acadêmicas frequentemente utilizam hardware proprietário ou de difícil aquisição no contexto brasileiro.

Como proposta de desenvolvimento, será empregado microcontroladores COTS (componentes comerciais prontos para uso, do inglês \textit{Commercial Off-The-Shelf}) , amplificadores operacionais, eletrodos de sEMG, softwares de simulação, programação e estruturas mecânicas fabricadas por impressão 3D, aliados a técnicas modernas de aprendizado de máquina para classificação de gestos.

